\section{Five Transports, and the One That Fails in Silence}
\label{sec:ch14-five-transports}
\index{subscriptions!the transport between router and subgraph|(}%

A subgraph entry in the composer's input can carry a \code{subscription:} block
with three keys, and between them they select how the router will reach that
service when a subscription needs it. Three transports, one of which has two
sub-flavours, so five arrangements in all: a WebSocket on either subprotocol,
server-sent events over a GET, and server-sent events over a POST. The
documentation for the file gives the keys as
follows~\autocite{cosmodocs2026compose}:

\begin{minted}[fontsize=\footnotesize,breaklines]{yaml}
    subscription:
      url: http://localhost:4001/sse #Optional, defaults to routing_url
      protocol: sse # Optional, defaults to ws (websockets)
      websocket_subprotocol: graphql-ws # Optional, defaults to auto.
      # Available options are auto, graphql-ws, graphql-transport-ws
\end{minted}

Since the last section ended on that third key being spelled differently in the
code, the honest way to find out what any of this does is to compose the real
eight subgraphs with the reviews entry varied, and read the block wgc wrote.
All seven variations below compose. None of them produces a message of any
kind.

\begin{minted}[fontsize=\scriptsize,breaklines]{text}
reviews subscription block                                 protocol      subprotocol
absent                                                     ..._WS        ..._AUTO
protocol: ws, websocketSubprotocol: graphql-transport-ws   ..._WS        ..._TRANSPORT_WS
protocol: ws, websocket_subprotocol: graphql-transport-ws  ..._WS        ..._AUTO
protocol: ws, websocketSubprotocol: totally-not-a-protocol ..._WS        (absent)
protocol: sse                                              ..._SSE       ..._AUTO
protocol: sse_post                                         ..._SSE_POST  ..._AUTO
protocol: carrier-pigeon                                   (absent)      ..._AUTO
\end{minted}

\index{wgc@\code{wgc}!what \code{router compose} does not validate}%
Row three is the documented spelling, and it leaves the setting at its default.
Row four is the right key with a value nothing implements, and it is worse:
rather than falling back, the composer drops the member from the config
entirely, so the result is not a default but an absence. Row seven does the same
thing to the protocol itself, and a subscription block with no protocol in it is
not a state any documentation describes.

None of this is a mystery once you read the code. The compose command reads the
camelCase key off the subscription block and nothing else, and wgc ships a
validator that rejects exactly the values in rows four and seven, by name, with
a good message. Every \code{subgraph} and \code{monograph} command calls it.
\code{router compose} does not.

That is the shape of the thing worth remembering rather than the specific
spelling, which some release will fix. The router-only path through this
toolchain, the one chapters~\ref{ch:composition} and~\ref{ch:enter-the-router}
chose deliberately because it needs no account and no control plane, is a
second-class path in the CLI, and it skips checks the registry path performs.
All seven rows are cases in \code{scripts/realtime-cases.mjs}, so a release that
fixes any of them fails the gate and gets read rather than passing quietly.

\subsection*{What each transport does to a HotChocolate subgraph}

\index{sse@\code{sse}!versus \code{sse\_post}}%
The two spellings in rows five and six look like the same thing and are not,
and this is where a wrong guess costs a working feature.

\code{ws} works, on either subprotocol, and the last section showed the
handshake for both.

\code{sse} does not work, at all, against HotChocolate 16.6.0. The router sends
this:

\begin{minted}[fontsize=\footnotesize,breaklines]{text}
GET /graphql HTTP/1.1
Host: host.docker.internal:5105
Accept: text/event-stream
Cache-Control: no-cache
\end{minted}

and gets this:

\begin{minted}[fontsize=\footnotesize,breaklines]{text}
HTTP/1.1 301 Moved Permanently
Location: http://host.docker.internal:5105/graphql/
\end{minted}

\index{MapGraphQL@\code{MapGraphQL}!the redirect to the IDE}%
A GET with no query on \code{MapGraphQL}'s route is a request for the GraphQL
IDE, which lives one slash further along, so HotChocolate redirects. The router
does not follow redirects. What reaches the client is:

\begin{minted}[fontsize=\footnotesize,breaklines]{text}
event: next
data: {"errors":[{"message":"Internal server error"}],"data":null}
\end{minted}

and then the stream closes. At the router's committed \code{log\_level: info}
nothing is logged about any of it, and the subgraph logs nothing at any level.
Turn the router up to \code{debug} and there is exactly one line:

\begin{minted}[fontsize=\footnotesize,breaklines]{text}
DEBUG abstractlogger@v0.0.4/zap.go:61 sseTransport.Subscribe
  {"endpoint": "http://localhost:5105/graphql", "method": "GET"}
\end{minted}

which says nothing about the redirect or the failure and does name the one
thing that matters, \code{"method": "GET"}. That is the whole diagnostic trail:
a one-word change in a configuration file, a composition that succeeds, an
error message naming nothing, and one debug line whose significance you have to
already know. I found the cause with a proxy in the middle and only then went
back and recognised the line.

\code{sse\_post} works:

\begin{minted}[fontsize=\footnotesize,breaklines]{text}
POST /graphql HTTP/1.1
Accept: text/event-stream
Cache-Control: no-cache
Content-Type: application/json

HTTP/1.1 200 OK
Content-Type: text/event-stream; charset=utf-8
Cache-Control: no-cache
\end{minted}

which is the request chapter~\ref{ch:schema-design-that-survives-change} made
by hand when it first exercised the subscription against one service. So that
transport is available, and of the two names, the one that sounds like it is
the one that is not.

\index{transports!choosing between them}%
Given that, why would anyone leave \code{ws}? Because a WebSocket is a
stateful connection to a subgraph, and a subgraph that holds thousands of them
is one you cannot scale to zero, cannot run on a request-scoped platform, and
have to drain carefully on deploy. \code{sse\_post} makes the router-to-subgraph
hop an ordinary HTTP request with a long response, which most infrastructure
already understands. Mosaic stays on \code{ws} because seven services on one
machine gain nothing from the swap, and because leaving it there is what makes
the next section's alternative worth comparing against.

\medskip
That alternative is the interesting one, and it takes the connection away from
the subgraph entirely.

\index{subscriptions!the transport between router and subgraph|)}%
